\chapter{Methodology: Job-Centric MARL Framework}
\label{chap:methodology}

Section~\ref{sec:gap} systematically filtered existing approaches to Dynamic Flexible Job-Shop Scheduling Problems (FJSP) through the requirement lens, excluding PDR-assisted methods (rule-based decision logic), single-agent DRL (scalability bottlenecks), and static E2E MARL (closed-world assumptions), ultimately identifying dynamic end-to-end multi-agent reinforcement learning as the most requirement-aligned approach class. However, the analysis revealed that current dynamic E2E MARL implementations introduce severe limitations that prevent them from addressing realistic FJSP characteristics under genuine stochastic conditions.

This chapter presents the job-centric MARL framework developed to bridge those gaps. The framework adapts MASA-QMIX (Chapter~\ref{chap:theoretical_background}) to enable operation in dynamic scheduling environments. Section~\ref{sec:overview} formalizes the problem as a job-centric Dec-POMDP. Section~\ref{sec:algorithm_selection} justifies MASA-QMIX selection based on algorithmic requirements. Sections~\ref{sec:scalability_adaptation} through~\ref{sec:action_masking} present seven architectural adaptations, each addressing a specific literature limitation through targeted modifications. For each adaptation, we present: (1) the observed limitation pattern, (2) root cause analysis, and (3) the solution mechanism. Finally, Sections~\ref{sec:architecture} through~\ref{sec:algorithm} provide implementation details.


\section{Overview: Job-Centric MARL Framework for Dynamic FJSP}
\label{sec:overview}

This research addresses a dual-resource Flexible Job-Shop Scheduling Problem (FJSP) under stochastic job arrivals, where production efficiency depends on coordinated real-time allocation of both machines and operators in a dynamic environment.

\subsection{Problem Context}

The job shop system consists of \( M \) machines, \( K \) operators, and a continuously evolving set of heterogeneous jobs \( J \). Each job \( j \in J \) comprises a sequence of operations \( O_j = \{o_{j,1}, o_{j,2}, \ldots, o_{j,n_j}\} \), where each operation must be executed on a compatible machine-operator pair before the next operation can begin.

\textbf{Key Problem Characteristics:}
\begin{itemize}
\item \textbf{Dual-Resource Allocation:}
Each operation requires simultaneous assignment of both a qualified machine and a qualified operator. The scheduling decision must select a feasible (machine, operator) pair from compatible resource combinations, where compatibility is determined by machine capability and operator skill qualifications.

\item \textbf{Stochastic Job Arrivals:}
New jobs enter the system during execution following a stochastic Poisson process. Unlike static scheduling scenarios, the scheduler operates without knowledge of future arrivals and must adapt policies in real-time as the workload evolves.

\item \textbf{Dynamic Agent Population:}
Each arriving job instantiates a new autonomous agent \( i_{\text{new}} \) that immediately participates in the learned cooperative policy without restarting the simulation or retraining the model. The number of active agents \( n(e) \) varies continuously throughout execution as jobs arrive and complete.

\item \textbf{Heterogeneous Jobs:}
Jobs are generated at runtime with randomly sampled operation sequences from an operation type pool. Jobs differ in operation types, sequence lengths, processing time requirements, and machine-operator qualification constraints. Processing times for the same operation may vary depending on which machine-operator pair executes it, reflecting real-world skill and capability differences.

\item \textbf{Operational Constraints:}
The scheduling environment operates under the following constraints:
\begin{itemize}
    \item Machine disjunction: Each machine processes at most one operation at a time,
    \item Operator disjunction: Each operator can be assigned to at most one machine at any given time,
    \item Precedence: Operations within a job must execute in sequential order,
    \item Non-preemption: Operations cannot be interrupted once started,
    \item Capability: Operations can only execute on (machine, operator) pairs where both the machine is capable and the operator is qualified,
    \item Deferrals: Jobs unable to find a feasible (machine, operator) pair wait until the next decision epoch.
\end{itemize}
\end{itemize}

\textbf{Optimization Objectives:}
The scheduling policy aims to minimize job waiting times (average flow time), minimize total makespan, maximize resource utilization and load balance, and maximize system throughput (jobs completed per time unit).

\subsection{Job-Centric MARL Formulation}

This work formulates scheduling as a Multi-Agent Reinforcement Learning (MARL) problem where \textbf{each job acts as an autonomous agent} responsible for selecting its own resource allocations. The problem follows the Decentralized Partially Observable Markov Decision Process (Dec-POMDP) framework established in Section~\ref{sec:problem_formulation}, defined by the tuple \( \langle \mathcal{N}, \mathcal{S}, \{\mathcal{A}^i\}, \mathcal{T}, \{\mathcal{O}^i\}, \mathcal{Z}, \mathcal{R}, \gamma \rangle \), where:

\begin{itemize}
    \item \( \mathcal{N} = \{1, 2, \ldots, n(e)\} \): Set of agents (active jobs) at decision epoch \( e \), where \( n(e) \) varies dynamically,
    \item \( \mathcal{S} \): Global state space encoding system-wide metrics (job counts, resource utilization, cumulative wait times),
    \item \( \mathcal{A}^i \subseteq \{1, \ldots, M\} \): Action space for agent \( i \), representing machine selection,
    \item \( \mathcal{O}^i \): Observation space for agent \( i \) (job-specific information: current operation type, remaining operations, wait time, available machines),
    \item \( \mathcal{R} \): Shared reward function balancing completion rate, wait time, throughput, and load balance,
    \item \( \gamma \): Discount factor for future rewards.
\end{itemize}

Agents learn cooperative policies \( \boldsymbol{\pi} = (\pi^1, \ldots, \pi^n) \) through shared rewards without requiring centralized coordination during execution. This job-centric design naturally handles dynamic agent populations and enables decentralized decision-making under partial observability—essential properties for real-time scheduling in stochastic production environments.

The objective is to maximize expected cumulative discounted reward:
\[
J(\boldsymbol{\pi}) = \mathbb{E}_{\tau \sim \boldsymbol{\pi}} \left[ \sum_{t=0}^{T} \gamma^t r_t \right]
\]

State representation, observation vectors, action spaces, and reward function components are detailed in the experimental setup (Chapter~\ref{chap:experiments}).


\section{Algorithm Selection: Why MASA-QMIX?}
\label{sec:algorithm_selection}

Having identified the problem formulation and requirements, the next question is: \textit{which MARL algorithm can support the necessary adaptations while maintaining computational efficiency and learning effectiveness?}

Chapter~\ref{chap:theoretical_background} reviewed MARL paradigms for cooperative multi-agent problems, categorizing algorithms into actor-critic methods (MADDPG, COMA, MAPPO) and value-based methods (VDN, QMIX, QTRAN). For dynamic FJSP, the problem characteristics impose several algorithmic requirements:

\begin{enumerate}
\item \textbf{Centralized Training with Decentralized Execution (CTDE):} Jobs must coordinate cooperatively to optimize system-wide objectives (minimize makespan, maximize throughput) while making independent resource selection decisions during execution without inter-agent communication.

\item \textbf{Discrete Action Spaces:} Resource selection is inherently discrete—agents choose from \( M \) machines. Value-based methods naturally handle discrete actions through Q-value maximization; actor-critic methods require categorical action distributions with additional complexity.

\item \textbf{Off-Policy Learning:} Stochastic arrivals create highly variable episode structures. Off-policy algorithms can learn from stored experience replay, reusing diverse episodes to improve sample efficiency.

\item \textbf{Computational Efficiency:} FJSP is NP-hard with exponential state-action spaces. The algorithm must achieve coordination through efficient value decomposition rather than expensive architectural additions (graph embeddings, attention mechanisms).

\item \textbf{Sufficient Cooperative Expressiveness:} The algorithm must capture coordination patterns beyond simple additive utility (VDN's limitation), requiring non-linear agent interaction modeling.

\item \textbf{Proven FJSP Applicability:} Algorithmic choices benefit from prior validation in scheduling domains.
\end{enumerate}

\textbf{MASA-QMIX Selection Rationale:}

QMIX satisfies all requirements through its core design (Section~\ref{sec:masa_qmix}):

\begin{itemize}
\item \textbf{CTDE via Monotonic Value Decomposition:} QMIX's mixing network enforces monotonicity (\( \partial Q_{\text{tot}} / \partial Q_i \geq 0 \)), guaranteeing the Individual-Global-Max (IGM) property: \( \arg\max_{\mathbf{a}} Q_{\text{tot}} = (\arg\max_{a^1} Q_1, \ldots, \arg\max_{a^n} Q_n) \). Each job independently selects its best action (decentralized execution) while the joint action remains globally optimal (centralized training objective).

\item \textbf{Discrete Action Spaces:} Q-learning foundation natively handles discrete actions through \( \varepsilon \)-greedy exploration and \( \arg\max \) action selection.

\item \textbf{Off-Policy Training:} Experience replay enables sample reuse across episodes with different stochastic arrival patterns.

\item \textbf{Computational Efficiency:} Simple feedforward networks (agent Q-networks, hypernetwork-generated mixer weights) without graph convolutions, attention layers, or communication protocols.

\item \textbf{Sufficient Expressiveness:} Non-linear mixing enables complex agent interactions while maintaining IGM guarantees.

\item \textbf{Scheduling Domain Validation:} Successfully applied to various scheduling problems, demonstrating effectiveness in discrete resource allocation tasks.
\end{itemize}

Additionally, QMIX's \textbf{parameter sharing} mechanism is critical: all agents use a single Q-network with shared parameters \( \theta \), enabling zero-shot generalization when new jobs arrive stochastically. A newly instantiated agent with \( h_0 = \mathbf{0} \) immediately uses the current policy without retraining—essential for continuous operation under dynamic arrivals.

However, standard MASA-QMIX implementations assume: (1) uniform timesteps \( t = 1, 2, 3, \ldots \), (2) fixed agent count \( n \) throughout episodes, (3) all agents acting at every timestep, and (4) independent action spaces without resource conflicts. Dynamic FJSP violates all four assumptions, requiring the architectural adaptations presented in subsequent sections.


\section{Adaptation 1 — Job-Centric Agent Design for Scalability}
\label{sec:scalability_adaptation}

\subsection{Observed Limitation in Literature}

Existing MARL-FJSP implementations exhibit severe scalability degradation as system workload increases. When workload grows (more concurrent jobs, higher arrival rates, longer operation sequences), system performance deteriorates: training becomes unstable, policies fail to generalize, and computational costs explode. Researchers respond by constraining problem sizes (e.g., maximum 10-15 concurrent jobs), simplifying action spaces (ignoring operator constraints, reducing machine options), or limiting job complexity (short operation sequences, homogeneous job types).

\subsection{Root Cause Analysis}

The scalability bottleneck stems from fixed agent populations in machine-centric architectures. When machines act as agents, the number of agents \( n \) remains constant regardless of workload—5 machines yield 5 agents whether the system processes 3 jobs or 30 jobs. As workload increases, each machine-agent must track exponentially more jobs in its observation space, reason about increasingly complex global states, and coordinate with other machines over growing action spaces. Furthermore, machine-agents exhibit low system awareness: machines observe all jobs but lack inherent knowledge of job-specific constraints (operation sequences, precedence requirements, due dates), forcing them to learn these relationships indirectly through trial and error.

\subsection{Proposed Solution: Job-Centric Agent Architecture}

This work adopts a job-centric architecture where \textbf{each job acts as an autonomous agent}. In the Dec-POMDP formulation (Section~\ref{sec:overview}), agent set \( \mathcal{N} = \{1, 2, \ldots, n(e)\} \) now corresponds to active jobs at decision epoch \( e \), not machines. This design choice addresses scalability through three mechanisms:

\begin{enumerate}
\item \textbf{Dynamic Agent Population:} Agent count \( n(e) \) varies with workload—when a job arrives stochastically at \( t_{\text{sim}} = 5.2 \), a new agent instantiates; when a job completes all operations, its agent exits. This direct coupling between workload and agent population eliminates fixed-size bottlenecks: the system naturally scales from 3 agents (low workload) to 30 agents (high workload) without architectural modification.

\item \textbf{High System Awareness:} Each job-agent maintains complete knowledge of its own state (current operation type, remaining operation sequence, accumulated wait time, resource requirements, precedence constraints). This role alignment—agents represent entities making scheduling decisions rather than passive resources being scheduled—simplifies observation construction and improves decision quality. Jobs inherently know what they need; machines must learn to infer it.

\item \textbf{Reduced Per-Agent Complexity:} Job-agents possess inherent self-knowledge—each job already knows its operation sequence, current operation requirements, precedence constraints, and processing characteristics without needing to learn or track them. This self-awareness means job-agents only need to observe system-wide resource states (which machines/operators are available, system workload snapshot) to make informed decisions. In contrast, machine-agents lack job-specific context: a machine observes that jobs exist but cannot inherently know each job's operation sequence length, which operation index it is on, which machines can perform its current operation, or its precedence dependencies. This information must either be explicitly provided (expanding observation dimensionality) or learned indirectly through experience (increasing learning complexity). The job-centric design eliminates this overhead—agents track only external system state rather than relearning their own characteristics, keeping per-agent observation complexity bounded by resource availability \( M + K \) instead of scaling with dynamic job population \( n(e) \).
\end{enumerate}


\section{Adaptation 2 — Runtime Job Generation for True Stochasticity}
\label{sec:stochastic_adaptation}

\subsection{Observed Limitation in Literature}

While recent work enables stochastic job arrivals in dynamic FJSP, implementations introduce severe simplifications that undermine real-world applicability. Four limitation patterns emerge:

\textbf{Pattern 1 — Controlled Stochasticity:} Studies enable Poisson arrivals but impose hard constraints on concurrent job counts (e.g., "maximum 10 jobs"). This transforms genuinely stochastic problems into static, controlled scenarios where policies learn to operate within predefined boundaries rather than adapting to uncontrolled, variable workloads.

\textbf{Pattern 2 — Homogenized Job Types:} Job diversity is severely restricted through single job types or small predefined template libraries (3-5 types). Agents optimize for template recognition ("when Job Type A arrives, allocate Machine 3 first") instead of learning generalizable decision heuristics based on system conditions.

\textbf{Pattern 3 — Processing Time Uniformity:} When job type diversity increases, processing time variability is systematically eliminated—all machines process operations in identical times. This neutralizes the routing optimization challenge central to FJSP, collapsing decisions from "which machine optimizes processing time?" to merely "which machine is available?"

\textbf{Pattern 4 — Action Space Coarsening:} Implementations apply aggressive filtering (availability-based, heuristic pre-selection, binary abstraction), reducing decision options to 2-3 discrete choices. Agents learn simplified dispatch rules instead of developing policies that reason over machine capabilities and coordination trade-offs.

\subsection{Root Cause Analysis}

The simplification stems from two interdependent factors: (1) \textbf{fixed agent architectures} cannot accommodate arbitrary agent counts (machine-agents require knowing maximum job count for observation sizing), and (2) \textbf{predefined job libraries} simplify training by limiting state-space exploration (agents encounter only familiar job types during training and deployment). Researchers constrain arrival rates, homogenize job characteristics, and coarsen action spaces to maintain training stability within these architectural limitations.

\subsection{Proposed Solution: Runtime Heterogeneous Job Generation}

This framework eliminates predefined job templates and arrival constraints through Poisson arrivals with runtime job construction:

\begin{enumerate}
\item \textbf{Unconstrained Poisson Arrivals:} Jobs arrive following Poisson process with rate \( \lambda \), with inter-arrival times sampled from \( \text{Exp}(1/\lambda) \). No episode-level constraints on concurrent job counts—system handles arbitrary arrival rates through job-centric scalability (Section~\ref{sec:scalability_adaptation}).

\item \textbf{Runtime Job Generation:} Each arriving job is constructed at arrival time by randomly sampling: operation types from operation pool (e.g., Op0--Op8), operation sequence length \( n_{\text{ops}} \sim \text{Uniform}(3, 8) \), and processing times from operation-machine lookup tables with stochastic variability reflecting machine capability differences. Jobs are runtime-generated entities with heterogeneous characteristics, not predefined templates.

\item \textbf{Heterogeneous Processing Times:} Processing time tables encode machine capability differences: newer equipment processes faster, high-precision machines handle complex operations more efficiently. This variability forces agents to learn \textit{capability-aware routing}—selecting machines that minimize processing time while balancing system-wide coordination.

\item \textbf{Full Action Space Exploration:} No action space coarsening or heuristic filtering. Agents explore full machine selection space \( \mathcal{A}^i \subseteq \{1, \ldots, M\} \) through \( \varepsilon \)-greedy policy, discovering nuanced routing strategies ("prefer Machine X for Operation Type Y under high utilization") through unconstrained exploration.

\item \textbf{Pattern-Neutral Learning:} Without predefined job templates, agents cannot memorize job-type-specific patterns. They must extract transferable situational patterns based on \textit{system conditions} (high workload \( \rightarrow \) defer; scarce machine \( \rightarrow \) anticipate conflicts) that generalize across arbitrary job types.
\end{enumerate}


\section{Adaptation 3 — Event-Driven Decision Architecture}
\label{sec:temporal_adaptation}

\subsection{Observed Limitation in Literature}

MARL training environments for FJSP predominantly employ fixed timesteps (\( t = 1, 2, 3, \ldots \)) where all agents make decisions synchronously at regular intervals. This temporal discretization creates three artifacts: (1) operations completing at irregular real times (e.g., \( t_{\text{sim}} = 2.37 \)) must wait until the next timestep (e.g., \( t = 3 \)) to request new resource allocations, introducing artificial delays; (2) agents forced to act at timesteps when no meaningful state changes occurred waste computation on redundant decisions; (3) policies learn coordination based on discretized time rather than true event-driven dynamics, degrading real-world transferability.

\subsection{Root Cause Analysis}

Fixed timesteps originate from episodic RL training paradigms designed for turn-based or uniformly-timed environments (e.g., Atari games, robotic control at fixed control frequencies). Researchers adopt this structure for FJSP implementations because: (1) \textbf{computational simplicity}—synchronous timesteps align with standard RL frameworks (Gym API, batch processing), avoiding the implementation complexity of asynchronous event handling and variable-length decision sequences, and (2) \textbf{compatibility with controlled scenarios}—when job arrivals are predetermined (Pattern 1 from Section~\ref{sec:stochastic_adaptation}) and processing times are uniform (Pattern 3), event timing becomes predictable, allowing timestep intervals to be experimentally tuned to approximate event occurrence (e.g., setting \( \Delta t = 1.0 \) to align with average operation completion times). 

However, this alignment collapses under true stochasticity with variable processing times: when jobs arrive following Poisson processes (unpredictable timing) and processing durations vary by machine capability (Operation A takes 3.2 time units on Machine 1 but 5.7 units on Machine 2), event occurrence times become impossible to predict or align with fixed intervals. Synchronous timesteps then force artificial delays (events at \( t_{\text{sim}} = 4.73 \) wait until \( t = 5 \)) and redundant decisions (agents act at \( t = 3 \) despite no state changes since \( t = 2 \)), distorting learning dynamics. FJSP is fundamentally event-driven—scheduling decisions arise when \textit{events} occur (job arrivals, operation completions), not at regular intervals.

\subsection{Proposed Solution: Discrete-Event Simulation Integration}

This work integrates MASA-QMIX with SimPy, a discrete-event simulation framework, enabling event-driven decision epochs:

\begin{enumerate}
\item \textbf{Continuous Simulation Time:} Time advances as \( t_{\text{sim}} \in \mathbb{R}^+ \) in continuous values (e.g., \( t_{\text{sim}} = 2.37, 5.91, 8.03, \ldots \)) corresponding to actual event occurrence times, not discretized steps.

\item \textbf{Event-Triggered Decision Epochs:} Agents make decisions only when events require them: job arrivals (need initial resource allocation) or operation completions (need next operation allocation). Between events, simulation advances without agent interaction, eliminating wasted computation.

\item \textbf{Asynchronous Agent Activation:} At each decision point (e.g., event \( e_i \) occurring at \( t_{\text{sim}} = 5.7 \)), only agents requiring decisions participate—jobs currently processing operations remain inactive. Decision batch size \( |\mathcal{B}_{e_i}| \) varies dynamically (e.g., 2 agents at event \( e_1 \), 5 agents at event \( e_2 \)) based on event timing, reflecting true system asynchrony.
\end{enumerate}


\section{Adaptation 4 — Parallel Decision-Making with Conflict Resolution}
\label{sec:conflict_adaptation}

\subsection{Observed Limitation in Literature}

Standard MARL implementations assume independent action spaces—each agent's action affects only its own state without direct interference. FJSP, however, involves explicit resource conflicts: multiple jobs can simultaneously require the same machine-operator pair. Existing work handles this implicitly through sequential action execution (Agent\textsubscript{1} selects first, state updates, Agent\textsubscript{2} observes updated state), creating order bias and first-come advantages that do not reflect true parallel decision-making.

\subsection{Root Cause Analysis}

The conflict resolution gap stems from misalignment between Dec-POMDP theory (assumes simultaneous joint actions) and implementation reality (actions execute sequentially in software). This sequential execution is not a design choice but a technical constraint: Python and PyTorch execute code statements sequentially—when multiple agents call \texttt{env.step(action)}, these calls must occur in some order (Agent\textsubscript{1} then Agent\textsubscript{2} then Agent\textsubscript{3}...). Each \texttt{step()} call updates environment state before the next agent observes, creating the order bias described above. Standard MASA-QMIX implementations do not address this because benchmark environments (e.g., StarCraft micromanagement) feature spatially separated agents with non-overlapping action targets—each unit controls different game entities, so sequential execution does not create resource conflicts.

\subsection{Proposed Solution: State Freezing with Explicit Conflict Resolution}

This framework implements true parallel decisions through state freezing and conflict detection mechanisms:

\begin{enumerate}
\item \textbf{State Freezing:} All agents observe identical frozen state \( s_{e_i} \) (state at event \( e_i \), capturing system configuration at that specific event time \( t_{\text{sim}} \)), select actions independently without observing others' choices.

\item \textbf{Conflict Detection:} After all agents in decision batch \( \mathcal{B}_{e_i} \) select actions, environment detects resource conflicts (multiple jobs requesting same machine-operator pair).

\item \textbf{Learned Conflict Resolution:} Conflicts are resolved via \( \varepsilon \)-greedy strategy: (1) during exploration (\( \varepsilon \) probability), random winner selected; (2) during exploitation, agent with highest Q-value wins. Losing agents re-queue for next decision epoch, storing conflict experience for learning.

\item \textbf{Proactive Conflict Avoidance:} Through experience replay, agents learn to estimate conflict probability for contested resources, developing strategies to defer or select alternative machines proactively.
\end{enumerate}

This adaptation ensures Dec-POMDP theoretical assumptions align with implementation reality, enabling agents to learn genuine coordination strategies rather than exploiting artificial execution order biases.


\section{Adaptation 5 — Asynchronous Agent Activation}
\label{sec:async_adaptation}

Standard MASA-QMIX assumes all \( n \) agents act at every timestep \( t \), regardless of whether they require decisions. In dynamic FJSP, this assumption wastes computation: jobs currently processing operations have no decisions to make until operation completion. This adaptation enables selective agent activation based on event triggers, improving computational efficiency and aligning with real-world asynchrony.

\textbf{Implementation:} At each decision epoch \( e_i \), only agents in decision batch \( \mathcal{B}_{e_i} \) participate—jobs that just arrived or completed operations. Jobs actively processing remain inactive. This selective activation is achieved through: (1) event-driven triggers (Section~\ref{sec:temporal_adaptation}) that identify which jobs require decisions, (2) dynamic batch construction containing only active agents, and (3) Q-network forward passes limited to \( |\mathcal{B}_{e_i}| \) agents instead of full population \( n(e) \).


\section{Adaptation 6 — Dynamic Population Handling}
\label{sec:dynamic_population}

Standard MASA-QMIX assumes fixed agent count \( n \) throughout episodes. Dynamic FJSP violates this: jobs arrive stochastically (increasing \( n \)) and complete (decreasing \( n \)). This adaptation enables seamless agent instantiation and removal without episode restarts.

\textbf{Implementation:} Parameter sharing (Section~\ref{sec:algorithm_selection}) enables zero-shot generalization: all agents use identical Q-network \( Q_{\theta}(o^i, a^i; h^i) \). When a job arrives at \( t_{\text{sim}} \), a new agent instantiates with initial hidden state \( h_0 = \mathbf{0} \) and immediately participates in policy \( \pi_{\theta} \) without retraining. When a job completes its final operation, the agent exits. QMIX's mixing network handles variable agent counts by zero-padding: if current population is \( n(e_i) = 7 \) but maximum observed was \( n_{\max} = 10 \), the mixer receives Q-values \( [Q_1, Q_2, \ldots, Q_7, 0, 0, 0] \), treating absent agents as contributing zero utility.


\section{Adaptation 7 — Action Masking for Dual-Resource Constraints}
\label{sec:action_masking}

FJSP's dual-resource constraints (machine capability, operator qualification) create variable action spaces: not all machines can perform all operations. Standard MASA-QMIX does not explicitly handle infeasible actions, allowing agents to select invalid machine-operator pairs that environment must reject, wasting exploration and destabilizing training.

\textbf{Implementation:} Action masking applies feasibility constraints before action selection: (1) for agent \( i \)'s current operation, identify compatible machines \( \mathcal{M}_{\text{capable}}(o_i) \) and qualified operators \( \mathcal{K}_{\text{qualified}}(o_i) \), (2) compute feasible action set \( \mathcal{A}_{\text{feasible}}^i = \{ m \mid m \in \mathcal{M}_{\text{capable}}, \exists k \in \mathcal{K}_{\text{qualified}} : \text{available}(m, k) \} \), (3) mask infeasible actions by setting Q-values to \( -\infty \): \( Q_i(o^i, a; h^i) \leftarrow -\infty \) for \( a \notin \mathcal{A}_{\text{feasible}}^i \), (4) apply \( \arg\max \) over masked Q-values, ensuring selected action is always feasible. This eliminates invalid exploration and accelerates convergence by focusing learning on feasible action subspace.


\section{Network Architecture}
\label{sec:architecture}

The job-centric MASA-QMIX framework consists of three neural components following the architecture detailed in Section~\ref{sec:masa_qmix}:

\textbf{Agent Q-Network:} Each job-agent uses a shared recurrent Q-network \( Q_{\theta}(o^i, a^i; h^i) \) with GRU hidden state to handle partial observability. Input dimension matches observation vector \( o^i \) (job state + system snapshot), output dimension equals action space size \( |\mathcal{A}^i| = M \) (number of machines).

\textbf{Mixing Network:} Combines individual Q-values into joint value \( Q_{\text{tot}}(\boldsymbol{\tau}, \mathbf{a}, s) \) via hypernetwork-generated weights. Hypernetwork inputs global state \( s \) and outputs mixing weights constrained to be non-negative (enforcing monotonicity for IGM property). The mixer handles variable agent counts through zero-padding as described in Section~\ref{sec:dynamic_population}.

\textbf{Target Networks:} Separate target networks \( Q_{\theta^-} \) and \( Q_{\text{tot}}^{\phi^-} \) stabilize training through periodic soft updates (\( \tau_{\text{polyak}} \)).

Network dimensions, activation functions, and hyperparameters are specified in experimental setup (Chapter~\ref{chap:experiments}).


\section{Training Strategy}
\label{sec:training_strategy}

Training integrates event-driven simulation with QMIX value decomposition:

\textbf{Experience Replay:} Decision epochs generate transitions \( (s_e, \mathbf{o}_e, \mathbf{a}_e, r_e, s_{e+1}, \mathbf{o}_{e+1}) \) stored in replay buffer. Variable-length episodes (different arrival patterns, stochastic durations) are handled through zero-padding and episode termination flags.

\textbf{TD Loss:} Standard QMIX TD-error (Section~\ref{sec:masa_qmix}) applies:
\[
\mathcal{L}(\theta, \phi) = \mathbb{E}_{(s, \mathbf{o}, \mathbf{a}, r, s', \mathbf{o}') \sim \mathcal{D}} \left[ \left( Q_{\text{tot}}(\boldsymbol{\tau}, \mathbf{a}, s; \phi) - y \right)^2 \right]
\]
where \( y = r + \gamma \max_{\mathbf{a}'} Q_{\text{tot}}^-(\boldsymbol{\tau}', \mathbf{a}', s'; \phi^-) \).

\textbf{Exploration Schedule:} \( \varepsilon \)-greedy exploration decays from \( \varepsilon_{\text{start}} = 1.0 \) to \( \varepsilon_{\text{end}} = 0.05 \) over training episodes.

\textbf{Batch Training:} Minibatches sampled from replay buffer update Q-networks and mixer via gradient descent. Target networks updated periodically (\( \tau_{\text{polyak}} = 0.005 \)).

Training hyperparameters (learning rates, buffer sizes, batch sizes, update frequencies) are detailed in Chapter~\ref{chap:experiments}.


\section{Software Architecture and Implementation}
\label{sec:implementation}

The system is implemented in Python using:
\begin{itemize}
\item \textbf{SimPy:} Discrete-event simulation framework for event-driven decision epochs
\item \textbf{PyTorch:} Neural network implementation (agent networks, mixer, training loop)
\item \textbf{Environment:} \texttt{environment.py} manages DES events, state/observation construction, and reward computation
\item \textbf{Agents:} \texttt{JobAgent} class represents jobs as SimPy processes with lifecycle management
\item \textbf{Resources:} \texttt{simpy.Resource} for machines, custom \texttt{OperatorPool} for operator management
\item \textbf{Replay Buffer:} \texttt{MARL/common/replay\_buffer.py} implements episodic storage with variable-length episode handling
\item \textbf{Training Loop:} \texttt{MARL/runner.py} orchestrates episode collection, batch sampling, and gradient updates
\end{itemize}

Modular design enables easy modification of arrival rates, job generation distributions, resource configurations, and reward function components for ablation studies and sensitivity analysis (Chapter~\ref{chap:experiments}).


\section{Algorithm: Job-Centric MASA-QMIX}
\label{sec:algorithm}

Algorithm~\ref{alg:jobcentric_masaqmix} presents the complete training and execution procedure, integrating all seven adaptations.

\begin{algorithm}[H]
\caption{Job-Centric MASA-QMIX for Dynamic FJSP}
\label{alg:jobcentric_masaqmix}

\KwIn{SimPy DES environment $E$, arrival rate $\lambda$, training epochs $M$}
\KwOut{Trained policy $\theta$ (agent network + mixer network)}

\textbf{Initialize:} Agent network $\theta_{\text{agent}}$, mixer network $\theta_{\text{mix}}$, target networks $\theta^-$, replay buffer $\mathcal{D}$, optimizer, exploration $\epsilon$\;

\For{$epoch = 1$ \KwTo $M$}{
    Reset environment $E$, $\mathcal{N}_{\text{active}} \leftarrow \emptyset$, start \texttt{TaskGenerator}($\lambda$)\;
    
    \tcp{Event-Driven Episode Loop}
    \While{not episode\_done}{
        \tcp{Stochastic Job Arrivals}
        \If{new job arrives $\sim$ Poisson($\lambda$)}{
            Create agent $i_{\text{new}}$ with $h_0 = \mathbf{0}$, add to $\mathcal{N}_{\text{active}}$\;
            Push decision request to \texttt{pending\_decisions}\;
        }
        
        \tcp{Advance Simulation Until Decision Epoch}
        Wait for \texttt{pending\_decisions} $\neq \emptyset$ (SimPy event loop)\;
        
        \tcp{Form Decision Batch (Asynchronous Activation)}
        $\mathcal{B}_e \leftarrow$ \texttt{pending\_decisions}.pop\_all()\;
        Freeze global state: $s_e \leftarrow$ get\_state($E$)\;
        
        \tcp{Parallel Action Selection with Conflict Detection}
        \ForEach{agent $i \in \mathcal{B}_e$}{
            $o^i \leftarrow$ get\_observation($i$)\;
            $\mathcal{M}^i \leftarrow$ compute\_mask($i, s_e$) \tcp{Three-condition feasibility}
            $Q^i \leftarrow$ AgentNetwork$(o^i, h_{t-1}^i; \theta_{\text{agent}})$\;
            $Q^i[\neg \mathcal{M}^i] \leftarrow -\infty$ \tcp{Mask infeasible actions}
            $a^i \leftarrow \epsilon$-greedy$(Q^i, \epsilon(t_{\text{sim}}))$\;
        }
        
        \tcp{Conflict Resolution}
        Conflicts $\leftarrow$ detect\_conflicts($\{a^i : i \in \mathcal{B}_e\}$)\;
        \ForEach{conflict $c$ in Conflicts}{
            Winners, Losers $\leftarrow$ resolve\_conflict($c$, mode=explore/exploit)\;
            Re-queue Losers to \texttt{pending\_decisions}\;
        }
        
        \tcp{Execute Actions and Advance Simulation}
        $(s', r) \leftarrow E$.execute\_actions(Winners)\;
        Store transition $(s_e, \{o^i, a^i\}, r, s')$ to buffer $\mathcal{D}$\;
        
        \tcp{Remove Completed Jobs}
        \ForEach{agent $i$ with all operations complete}{
            Remove $i$ from $\mathcal{N}_{\text{active}}$\;
        }
    }
    
    \tcp{Training Phase (Centralized Learning)}
    \If{$|\mathcal{D}| \geq$ batch\_size}{
        Sample batch of episodes from $\mathcal{D}$\;
        \ForEach{transition $(s, \{o^i, a^i\}, r, s')$ in batch}{
            $Q^i_{\text{eval}} \leftarrow$ AgentNetwork$(o^i, a^i; \theta_{\text{agent}})$ for all $i$\;
            $Q^i_{\text{target}} \leftarrow$ TargetNetwork$(o'^i; \theta^-)$ for all $i$\;
            $Q^i_{\text{target}}[\neg \mathcal{M}'^i] \leftarrow -\infty$\;
            
            $Q_{\text{tot}} \leftarrow$ MixingNetwork$([Q^1_{\text{eval}}, \ldots, Q^{n(e)}_{\text{eval}}], s; \theta_{\text{mix}})$\;
            $Q_{\text{target}} \leftarrow$ MixingNetwork$(\max_a Q^i_{\text{target}}], s'; \theta^-_{\text{mix}})$\;
            
            $y \leftarrow r + \gamma \cdot Q_{\text{target}}$\;
            $\mathcal{L} \leftarrow (Q_{\text{tot}} - y)^2$\;
            $\theta \leftarrow \theta - \alpha \cdot$ clip\_grad$(\nabla_\theta \mathcal{L})$\;
        }
        
        \tcp{Target Network Update}
        \If{train\_step $\mod C = 0$}{
            $\theta^- \leftarrow \theta$\;
        }
    }
    
    \tcp{Decay Exploration Based on SimPy Time}
    $\epsilon(t_{\text{sim}}) \leftarrow \max(\epsilon_{\text{end}}, \epsilon_{\text{start}} - \frac{t_{\text{sim}}}{T_{\text{anneal}}}(\epsilon_{\text{start}} - \epsilon_{\text{end}}))$\;
}

\Return{$\theta = (\theta_{\text{agent}}, \theta_{\text{mix}})$}\;

\end{algorithm}

\textbf{Key Algorithm Features:}
\begin{itemize}
\item \textbf{Line 6-8:} Stochastic job arrivals via Poisson process (Section~\ref{sec:stochastic_arrivals})
\item \textbf{Line 11:} Event-driven decision epochs (Section~\ref{sec:event_driven})
\item \textbf{Line 14:} Asynchronous agent activation—only decision-ready jobs participate (Section~\ref{sec:asynchronous_activation})
\item \textbf{Line 15:} State freezing for parallel decisions (Section~\ref{sec:conflict_resolution})
\item \textbf{Line 18-19:} Three-condition action masking (Section~\ref{sec:action_masking})
\item \textbf{Line 24-28:} Explicit conflict detection and resolution (Section~\ref{sec:conflict_resolution})
\item \textbf{Line 33-35:} Dynamic population—completed jobs exit seamlessly (Section~\ref{sec:dynamic_population})
\item \textbf{Line 52:} SimPy-time-based exploration decay (Section~\ref{sec:training})
\end{itemize}
